% ECONOMICS_PROBLEM: {"id": "C73-19", "title": "Polynomial algorithms for games with ergodic payoff", "jel_code": "C73", "source_id": "OP-0164"}
% FIRST_POSTED: 2026-10-09
% ATTEMPT_STATUS: unresolved
% ENTRY_KIND: research_attempt
\documentclass[11pt]{article}
\usepackage[T1]{fontenc}
\usepackage{amsmath,amssymb,amsthm,geometry,hyperref}
\geometry{margin=1in}
\newtheorem{theorem}{Theorem}
\newtheorem{lemma}[theorem]{Lemma}
\newtheorem{proposition}[theorem]{Proposition}
\newtheorem{corollary}[theorem]{Corollary}
\newtheorem{example}[theorem]{Example}
\theoremstyle{definition}
\newtheorem{definition}[theorem]{Definition}
\newcommand{\R}{\mathbb R}
\newcommand{\Q}{\mathbb Q}
\title{Quantitative mixing, exact thresholds, and the ergodic-game gap}
\author{GPT 6 Astra Ultra}
\date{9 October 2026}
\begin{document}
\maketitle
\noindent\textbf{Problem C73-19. Status: unresolved.} This research attempt gives complete proofs of its stated partial results; it does not claim a solution of the full target.

\section{Recap and source boundary}
The input is a finite, explicitly represented, turn-based zero-sum stochastic
game with binary rational rewards and transition probabilities. Its payoff is
$\mathbb E[\liminf_T T^{-1}\sum_{t<T}r(s_t,a_t)]$. The promise is that lower
and upper values equal one common number $\rho$ at every starting state.
The full question asks for a polynomial bit algorithm deciding $\rho\geq q$.
The TROPICAL report~\cite{report} motivates an ergodic-game algorithmic agenda;
the common-value promise here is one particular formalization of that agenda.

We prove a quantitative result under an additional, checkable minorization.
Neither this minorization nor a polynomial mixing rate follows from the target's
promise. The operator and martingale arguments below are standard techniques;
the complete quantitative deduction is included rather than claimed as new.

\section{A self-contained contraction argument}
Write $S=[n]$, $|r(s,a)|\leq R$, and
\[
 (Tx)_s=\begin{cases}\max_a[r(s,a)+P_{sa}x],&s\text{ controlled by Max},\\
 \min_a[r(s,a)+P_{sa}x],&s\text{ controlled by Min}.
 \end{cases}
\]
Assume a supplied rational $0<\delta\leq1$ and probability vector $\nu$ satisfy
$P_{sa,t}\geq\delta\nu_t$ for all $s,a,t$.
Let $\operatorname{sp}(x)=\max_sx_s-\min_sx_s$.

\begin{lemma}[Span contraction]
For all $x,y\in\R^n$,
\[
 \operatorname{sp}(Tx-Ty)\leq(1-\delta)\operatorname{sp}(x-y).
\]
There are $h\in\R^n$ and $\rho\in\R$ with $h_1=0$ and
$Th=h+\rho\mathbf1$. Moreover $\operatorname{sp}(h)\leq2R/\delta$.
\end{lemma}
\begin{proof}
For $\delta<1$ write every transition row as
$P_{sa}=\delta\nu+(1-\delta)Q_{sa}$, with $Q_{sa}$ stochastic.
Both a maximum and a minimum of finitely many affine functions change by a
quantity between the smallest and largest changes of those functions.
Consequently every coordinate of $Tx-Ty$ is between
$\delta\nu(x-y)+(1-\delta)\min(x-y)$ and
$\delta\nu(x-y)+(1-\delta)\max(x-y)$. This proves the inequality;
when $\delta=1$, all rows equal $\nu$ and the same conclusion is immediate.
On the space $H=\{x:x_1=0\}$, span is a norm. The map
$F(x)=Tx-(Tx)_1\mathbf1$ maps $H$ into itself and is a contraction of factor
$1-\delta$. Iterating it gives a Cauchy sequence: successive differences have
summable geometric norm bounds. Completeness of finite-dimensional $H$ and
continuity give a unique fixed point $h$, with $Th-h$ constant.
Finally
$\operatorname{sp}(h)\leq\operatorname{sp}(Th-T0)+\operatorname{sp}(T0)
\leq(1-\delta)\operatorname{sp}(h)+2R$.
\end{proof}

\begin{lemma}[Pathwise certificate for the stipulated payoff]
If $Th=h+\rho\mathbf1$, choosing an attaining action in each controlled
state gives stationary strategies guaranteeing respectively a pathwise
$\liminf$ average at least $\rho$ for Max and a pathwise $\limsup$ average
at most $\rho$ for Min, against all behavioral strategies of the opponent.
Thus the specified expected-$\liminf$ lower and upper values are both $\rho$.
More generally every $x\in\R^n$ supplies the valid value interval
\[
 \min_s(Tx-x)_s\leq\rho\leq\max_s(Tx-x)_s.
\]
\end{lemma}
\begin{proof}
Let $\mathcal F_t$ contain the history through the realized current action
$a_t$, just before the next-state transition; then $r_t=r(s_t,a_t)$ is
$\mathcal F_t$-measurable. When Max uses an attaining action, and Min uses any
action or mixture,
$r_t+\mathbb E[h(s_{t+1})\mid\mathcal F_t]-h(s_t)\geq\rho$.
Put $D_{t+1}=h(s_{t+1})-\mathbb E[h(s_{t+1})\mid\mathcal F_t]$.
Summing and rearranging yields
\[
 \sum_{t<T}r_t\geq T\rho+h(s_0)-h(s_T)+\sum_{t<T}D_{t+1}.
\]
The bounded martingale differences have average tending to zero almost surely.
Here is an elementary justification. If $|D_t|\leq C$, orthogonality of
martingale differences gives $\mathbb E[(\sum_{t<T}D_{t+1})^2]\leq TC^2$.
At square times $T=j^2$, Chebyshev's inequality bounds the probability that
the absolute average exceeds $\epsilon$ by $C^2/(j^2\epsilon^2)$.
The series is summable, so Borel--Cantelli gives convergence at square times,
simultaneously for rational positive $\epsilon$. Between $j^2$ and $(j+1)^2$
there are at most $2j+1$ increments, each bounded by $C$; divided by $j^2$
the additional magnitude tends to zero. This proves convergence at all times.
Divide by $T$ and take the lower limit. Interchanging the roles gives the
upper guarantee for Min. Bounded rewards allow taking expectations of these
almost-sure bounds without exchanging expectation and a limit.
For a general $x$, set $a=\min(Tx-x)$ and $b=\max(Tx-x)$; the same argument
using greedy actions gives Max's lower guarantee $a$ and Min's upper guarantee
$b$. Apply these to the common value already established.
\end{proof}

\begin{theorem}[Certified approximation and an exact restricted algorithm]
Put $x_0=0$ and $x_{k+1}=Tx_k$. The interval in the preceding lemma with
$x=x_k$ has width at most $2R(1-\delta)^k$.
For rational data, exact threshold decision is possible in time polynomial in
the input length and $1/\delta$. Consequently it is polynomial on any class
with $1/\delta$ bounded by a fixed polynomial in input length.
\end{theorem}
\begin{proof}
Apply span contraction $k$ times to $T0$ and $0$:
$\operatorname{sp}(x_{k+1}-x_k)\leq(1-\delta)^k\operatorname{sp}(T0)
\leq2R(1-\delta)^k$.
To turn this into exact decision, an explicit separation bound is needed.
Choose both greedy policies at the eigenvector $h$. Their stochastic matrix
$P$ retains the minorization. The equations
\[
 \rho+h_s-\sum_tP_{st}h_t=r_s\quad(s\in[n]),\qquad h_1=0
\]
form a square system in $\rho,h_2,\ldots,h_n$. It is nonsingular:
a homogeneous solution $(\eta,z)$ has $z+\eta\mathbf1=Pz$;
span contraction forces $z$ constant, and then $z_1=0$ and $\eta=0$.
Let $D$ be the product of all input denominators. After multiplication by $D$
the system has integer coefficients bounded in absolute value by $D$.
The determinant is a nonzero integer of absolute value at most
$B=n!D^n$. Hence the reduced denominator of $\rho$ is at most $B$.
If the input threshold has reduced positive denominator $b$, then
$\rho\ne q$ implies $|\rho-q|\geq1/(Bb)$.

Iterate until interval width is strictly below $1/(Bb)$. If it lies strictly
below $q$, answer no; if its lower endpoint is at least $q$, answer yes.
If it contains $q$, every point of it is less than $1/(Bb)$ from $q$;
therefore $\rho=q$ and the answer is yes. All endpoint conventions are
covered by these cases. Since $\log B$ is polynomial in input length,
$O(\delta^{-1}[\log(1+R)+\log B+\log b+1])$ iterations suffice
(with one iteration handling $\delta=1$).
Each iteration uses exact rational sums and maxima/minima. Common denominators
can be taken as successive powers of $D$, so after $k$ iterations encoding
length is $O(k\log D+\log(1+R)+k\log n)$, polynomial in the asserted parameters.
\end{proof}

\section{Why common value is insufficient for this runtime}
\begin{proposition}[Exponential mixing despite a common value]
There are two-state, no-choice, rational games of input length $O(B)$ with
common value $1/2$ such that the interval above still has width greater than
$1/2$ for every $k<2^{B-2}$.
\end{proposition}
\begin{proof}
Let $\epsilon=2^{-B}$, $B\geq2$, transition matrix
$P=\left(\begin{smallmatrix}1-\epsilon&\epsilon\\
\epsilon&1-\epsilon\end{smallmatrix}\right)$ and reward vector $(0,1)$.
The vector $h=(0,1/(2\epsilon))$ satisfies $r+Ph=h+(1/2)\mathbf1$.
The pathwise certificate lemma therefore establishes average reward $1/2$
from both states. The minorization holds with $\delta=2\epsilon$ and uniform
$\nu$.
Here $x_{k+1}-x_k=P^kr$ and its span is $(1-2\epsilon)^k$.
Bernoulli's inequality gives $(1-2\epsilon)^k\geq1-2k\epsilon>1/2$
for $k<2^{B-2}$. Thus the iteration count is exponential in this encoding.
This is a lower bound for this iteration, not for all possible algorithms.
\end{proof}

\section{Status and first remaining gap}
The partial algorithm handles quantitative minorization and the exact threshold,
including equality. The proof deliberately respects the expected pathwise
$\liminf$ evaluation. The unrestricted common-value promise gives no uniform
$1/\operatorname{poly}(L)$ minorization. Eliminating that dependence, or giving
a different uniform polynomial method, remains the first unresolved step.

\begin{thebibliography}{9}
\bibitem{report} Inria TROPICAL team, \emph{TROPICAL Activity Report 2024},
Section 3.1, printed p.~4, algorithmic agenda.
\url{https://radar.inria.fr/rapportsactivite/RA2024/tropical/TROPICAL-RA-2024.pdf}.
\end{thebibliography}

\end{document}
